\documentclass[aps,prmaterials,reprint,amsmath,amssymb,superscriptaddress,longbibliography]{revtex4-2}

\usepackage{graphicx}
\usepackage{bm}
\usepackage{booktabs}
\usepackage{siunitx}
\usepackage{xcolor}
\usepackage{hyperref}
\usepackage{microtype}

\newcommand{\Teff}{T_{\mathrm{eff}}}
\newcommand{\taufp}{\tau_{\mathrm{FP}}}
\newcommand{\etaperp}{\eta_{\mathrm{inter}}}
\newcommand{\kb}{k_{\mathrm B}}
\newcommand{\dd}{\mathrm d}

\begin{document}

\title{Dimensional rescue and tensorial hopping transport in three-dimensional atomic networks}

\author{J. P. Dadario Pereira}
\affiliation{Applied Physics Department, Gleb Wataghin Institute of Physics, State University of Campinas, Campinas, SP 13083-970, Brazil}
\author{Raphael Tromer}
\affiliation{Institute of Physics, University of Bras\'ilia, Bras\'ilia, Federal District 70910-900, Brazil}
\author{Luiz A. Ribeiro Junior}
\affiliation{Computational Materials Laboratory, LCCMat, Institute of Physics, University of Bras\'ilia, Bras\'ilia, Federal District 70910-900, Brazil}
\author{Douglas S. Galv\~ao}
\affiliation{Applied Physics Department, Gleb Wataghin Institute of Physics, State University of Campinas, Campinas, SP 13083-970, Brazil}

\date{\today}

\begin{abstract}
Hopping transport in disordered materials is controlled not only by local transition rates but also by the dimensionality and connectivity of the network on which those transitions occur. We develop a trajectory-resolved kinetic Monte Carlo framework that extends an atomistic two-dimensional hopping model to genuinely three-dimensional and chemically heterogeneous atomic networks. The formulation separates three ingredients that are often conflated: graph connectivity, kinetic accessibility, and local chemical energetics. In a controlled four-layer carbon network, vacancies progressively fragment strictly two-dimensional transport, whereas the activation of transverse links restores source--drain connectivity with remarkable efficiency. At a vacancy fraction of 0.25, the mean fraction of source sites connected to the drain increases from 0.275 in the strict-2D limit to 0.990 when transverse edges are available, while the probability that all source sites are connected rises from zero to 0.878. Once connectivity is restored, increasing the transverse hopping scale reduces the median reduced first-passage time by more than a factor of five, revealing a distinct crossover from topological rescue to kinetic acceleration. Periodic simulations on an atomistic BN crystal yield a full diffusion tensor whose slow principal component is strongly suppressed at weak transverse coupling; at a representative coupling of 0.05 the principal anisotropy is approximately seven. Finally, an empirical chemistry-aware model on W$_2$O$_6$ shows that local energetic contrast can redirect successful trajectories toward an O-rich backbone and change first-passage times by more than an order of magnitude even when eventual passage remains unity. These results establish dimensionality as an independent control parameter for hopping networks and provide a compact route from CIF structures to topology, trajectories, and tensorial transport observables.
\end{abstract}

\maketitle

\section{Introduction}

Charge motion in structurally complex materials is frequently governed by a competition between localization, thermal activation, and the existence of connected pathways across a disordered energy landscape. The conceptual foundations of this regime span Anderson localization, phonon-assisted impurity conduction, variable-range hopping, and percolative descriptions of disordered networks \cite{anderson1958,miller1960,ambegaokar1971,mott1979,shklovskii1984,abrahams1979,kramer1993,bassler1993}. In the strongly localized limit, the microscopic object that matters is no longer a Bloch state extending through the sample, but a sequence of stochastic transitions among spatially localized states. This change of description is especially important near connectivity thresholds, where a modest geometric or energetic perturbation can reorganize the small subset of links that carries long-range transport.

Kinetic Monte Carlo (KMC) provides a natural language for this problem because it converts a set of local rates into explicit real-time trajectories. Continuous-time Monte Carlo algorithms have therefore become standard tools for activated dynamics \cite{bortz1975,andersen2019}. In electronic materials, atomistic and molecular KMC schemes have been used to connect morphology, site energies, transfer integrals, and charge-transfer rates to mobilities, current profiles, and residence-time distributions \cite{ruhle2011,vehoff2010,mesta2013,liu2018,balzer2021}. Recent work continues to emphasize that three-dimensional percolation and variable-range connectivity can qualitatively alter transport in disordered semiconductors \cite{liu2025}. These studies also make clear that ``three-dimensional KMC'' by itself is not a novelty; the physically useful question is which aspects of three-dimensionality are responsible for changes in transport and how those changes can be separated from energetic disorder.

Two-dimensional materials provide a particularly transparent setting in which to formulate that question. Graphene combines a simple atomic topology with exceptionally rich transport phenomenology \cite{novoselov2004,geim2007,castroneto2009,peres2010}. In clean samples, high mobility and long mean free paths motivate coherent or semiclassical descriptions \cite{morozov2008,bolotin2008,hwang2008}. Vacancies and chemical disorder, however, generate localized states and strong scattering \cite{pereira2006,ugeda2010,chen2009}, and sufficiently disordered graphene can enter regimes where hopping and percolation become useful effective descriptions \cite{lherbier2008,lherbier2012,liang2012}. Similar defect-mediated hopping has been observed in other two-dimensional systems, including MoS$_2$ \cite{qiu2013}. In this context, topology is not a secondary detail: disorder removes redundant pathways, and the transport response becomes controlled by a sparse backbone whose geometry is sample specific.

Layered crystals then raise an immediate extension of the two-dimensional problem. A weak interlayer channel may be kinetically slow yet topologically decisive, because it can connect two in-plane components that would otherwise be disjoint. The distinction between a weak link and an absent link is therefore singular near a percolation threshold. Hexagonal boron nitride (h-BN) offers a useful atomistic geometry for examining this issue because its structure is strongly layered and its electronic response is anisotropic \cite{laturia2018,jiang2018}. Defects in h-BN generate localized electronic states \cite{huang2012,sajid2018}, and tunnel experiments have directly resolved localized states and multistep percolative processes through BN barriers \cite{greenaway2018}. More generally, defect-assisted interlayer processes are relevant to van der Waals materials and heterostructures \cite{prasad2023}. Our use of BN below is deliberately geometric and phenomenological: it provides a realistic layered crystal on which to test tensorial hopping transport rather than a parameter-free prediction of intrinsic BN conductivity.

Dimensionality also interacts with chemical heterogeneity. Microscopic hopping rates can depend on site-energy differences, electronic couplings, reorganization energies, and orbital character \cite{ruhle2011,slater1954}. Fully first-principles parameterization is possible in favorable cases, but it is expensive because a large disordered structure may contain an enormous number of inequivalent sites and pairs \cite{setyawan2010,ruhle2011}. In many situations the more informative strategy is to separate robust topology-driven effects from uncertain microscopic scales. This is analogous to the long-standing use of phenomenological disorder distributions and reduced transport variables in hopping theory \cite{shklovskii1984,bassler1993}. A model can then ask which conclusions survive broad parameter changes before attempting a material-specific electronic-structure calibration.

Our previous work introduced a trajectory-resolved hopping framework for planar carbon networks, with atomic coordinates read directly from crystallographic files and transport generated by continuous-time KMC \cite{pereira2026}. That formulation was intentionally simple and emphasized how vacancies, strain, temperature, and external fields reshape successful carrier paths. Extending the same framework to three dimensions exposed two methodological issues that are easy to miss in an effective current--voltage analysis: a finite observation window can masquerade as localization, and a globally imposed gap-activation factor can generate an artificial voltage threshold. The present work removes those ambiguities and reformulates the model around graph reachability, first-passage statistics, and diffusion tensors. The Landauer normalization used previously as a diagnostic mapping of trajectory success is not used here to assign a bulk three-dimensional conductivity \cite{landauer1957}.

The resulting framework has three levels, summarized in Fig.~\ref{fig:overview}. First, a controlled four-layer carbon stack isolates the role of dimensionality by comparing a strictly two-dimensional graph with the same structure after transverse edges are enabled. Second, periodic KMC on a real BN crystal yields the full diffusion tensor and its principal axes. Third, a W$_2$O$_6$ network introduces a minimal chemical contrast through species-dependent site energies and pair-dependent hopping weights. Across these levels we show that dimensionality acts in two separable stages: weak transverse connectivity can rescue a fragmented network, whereas stronger transverse coupling accelerates paths that already exist. Chemical contrast then selects which subset of the connected network is dynamically active. This separation of topology, kinetics, and chemistry is the central result of the paper.

\begin{figure*}[t]
\centering
\includegraphics[width=0.98\textwidth]{../reproduction/generated/figure1_system_overview.png}
\caption{\textbf{Hierarchy of the three-dimensional hopping framework.} (a) Atomic coordinates are read from a CIF file, converted into a geometric graph, and propagated by continuous-time KMC to obtain trajectory-resolved observables. (b) A finite four-layer carbon device separates source--drain graph connectivity from the kinetic cost of transverse hopping. Random vacancies remove in-plane redundancy, whereas interlayer links controlled by $\etaperp$ provide bypasses. (c) Periodic KMC produces the diffusion tensor $\mathbf D$ and its principal components. (d) Chemical heterogeneity is introduced through site-energy offsets and pair-dependent hopping weights. The hierarchy is intentionally phenomenological: it isolates mechanisms before material-specific electronic-structure calibration.}
\label{fig:overview}
\end{figure*}

\section{Methods}

\subsection{Atomic graphs and transport coordinates}
Atomic structures are read directly from crystallographic information files (CIFs). Each atom defines a hopping site at Cartesian position $\mathbf r_i$. For finite-device calculations, sites $i$ and $j$ are connected whenever their Euclidean separation $r_{ij}=|\mathbf r_j-\mathbf r_i|$ is below a prescribed cutoff. Neighbor construction uses a spatial tree rather than an all-pairs search, allowing the graph to remain sparse for larger three-dimensional systems.

A transport direction is represented by a Cartesian unit vector $\hat{\mathbf n}$ and defines the scalar coordinate $u_i=\mathbf r_i\cdot\hat{\mathbf n}$. For crystallographic transport along a nonorthogonal cell direction, $\hat{\mathbf n}$ is chosen from the corresponding reciprocal (dual) basis vector, i.e., normal to the transverse lattice plane. This convention is invariant under periodic translations in the two transverse directions and avoids the spurious longitudinal displacement that results from using a direct lattice vector in a skew cell. The source and drain are finite slabs adjacent to the minimum and maximum $u_i$, respectively.

\subsection{Site energies and hopping rates}
An applied bias $V$ is represented by a linear electrostatic drop over the device length $L=u_{\max}-u_{\min}$,
\begin{equation} E_i(V)=\epsilon_i-eV\frac{u_i-u_{\min}}{L}, \end{equation}
where $\epsilon_i$ is a local phenomenological site-energy offset. The transition rate from $i$ to a connected site $j$ is a Miller--Abrahams form \cite{miller1960},
\begin{align}
\Gamma_{i\rightarrow j}&=\nu_0 w_{ij}e^{-2r_{ij}/\xi}A_{ij},\\
A_{ij}&=\begin{cases}e^{-(E_j-E_i)/(\kb T)}, & E_j>E_i,\\1, & E_j\leq E_i.\end{cases}
\end{align}
The stochastic waiting time is $\Delta t=-\ln s/\sum_j\Gamma_{i\rightarrow j}$, with the next site selected proportionally to its rate. The main results do not use a global multiplicative activation factor tied to a band gap; the Supplemental Material documents that ablation. We emphasize the reduced first-passage time $\nu_0\taufp$ whenever possible.

\subsection{Dimensional crossover and exact graph reachability}
For layered systems, an edge is classified as transverse when its displacement has a component along the layer normal exceeding a fixed geometric threshold. Their relative kinetic strength is controlled by $w_{ij}=1$ for intralayer edges and $w_{ij}=\etaperp$ for interlayer edges, with $\etaperp=0$ defining the strict-2D limit. To separate topology from dynamics, exact source--drain reachability is computed on the active graph before KMC. The source reachability fraction is $P_{\mathrm{reach}}=N_{\mathrm{source}\rightarrow\mathrm{drain}}/N_{\mathrm{source}}$.

\subsection{First-passage observables and censoring}
Finite-device trajectories terminate when the carrier reaches the drain, exceeds a maximum number of hops, or exceeds the observation time. The effective passage fraction $\Teff$ is the fraction of launched trajectories that reach the drain. For successful paths we store $\taufp$, hop count, accumulated path length, and tortuosity. A fixed observation window can incorrectly classify long but connected trajectories as nontransmitting. We therefore use adaptive windows in validation calculations and, for the production dimensional-crossover statistics, condition kinetic sampling on source sites known by exact graph analysis to reach the drain.

\subsection{Periodic diffusion tensor}
Bulk-like diffusion is obtained from unbiased KMC with periodic boundary conditions,
\begin{equation}D_{\alpha\beta}=\frac{\langle\Delta r_\alpha\Delta r_\beta\rangle}{2t}.\end{equation}
When nonzero site-energy offsets are present, walkers are initialized from the equilibrium site distribution $P_i\propto\exp(-\epsilon_i/\kb T)$. The symmetric tensor is diagonalized to obtain principal diffusivities $D_1\geq D_2\geq D_3$ and eigenvectors. Directional projections on crystallographic plane normals provide an independent check. An Einstein mobility tensor $\bm\mu=e\mathbf D/(\kb T)$ can also be reported where appropriate \cite{baranovskii1996}.

\subsection{Model systems and statistical protocols}
The controlled dimensional-crossover system is a four-layer carbon stack containing 192 sites before vacancies. Exact topology statistics use 600 independent vacancy realizations at each vacancy fraction from 0 to 0.40. Conditional kinetics are evaluated at $p_v=0.25$ for 20 disorder realizations and $\etaperp=0.01,0.02,0.05,0.10,0.20,$ and 1.0; each point uses 800 trajectories at 300 K and 0.05 V.

The tensor benchmark uses the supplied BN crystal replicated $2\times2\times2$ to 64 atoms. For each $\etaperp$, 12 independent calculations of 1800 walkers are propagated for $1.2\times10^{-10}$ s. The chemistry benchmark uses a 48-atom W$_{12}$O$_{36}$ cell. We sample $\Delta\epsilon=0,0.25,0.50,$ and 1.0 eV and $w_{\mathrm{homo}}=1,0.3,0.1,$ and 0.03, with 20 stochastic realizations of 1200 trajectories per point. These parameters are deliberately empirical; the calculation is a sensitivity analysis rather than a quantitative prediction for tungsten oxide.

\section{Results and Discussion}
\subsection{Dimensional rescue is a topology transition before it is a kinetic acceleration}
Figure~\ref{fig:rescue} isolates dimensionality in the controlled layered network. At $p_v=0.25$, the mean source reachability in the strict-2D graph is $0.275\pm0.008$ (standard error), while allowing transverse edges raises it to 0.990. None of the 600 strict-2D disorder realizations is fully source connected, whereas 87.8\% of the 3D-enabled realizations are. Because these numbers come from exact graph traversal, no choice of $\nu_0$, temperature, localization length, or observation time creates the effect: the third dimension is functioning as a redundancy channel in the percolation network.

The kinetic statistics reveal a second regime. Restricting injection to source sites already known to reach the drain gives $\Teff=1$ and zero final censoring. Nevertheless, the median $\nu_0\taufp$ decreases from 1262 at $\etaperp=0.01$ to 247 at $\etaperp=1.0$, while tortuosity falls from 66.3 to roughly 26--29. Thus the existence of a transverse connection repairs topology, while its strength determines the dynamical price of the rescued detour.

\begin{figure*}[t]\centering
\includegraphics[width=0.98\textwidth]{../reproduction/generated/figure2_dimensional_rescue.png}
\caption{\textbf{Dimensional rescue in a vacancy-fragmented layered network.} Exact graph reachability and conditional first-passage kinetics separate topological reconnection from kinetic acceleration. Error bars and shaded intervals quantify disorder ensembles.}\label{fig:rescue}\end{figure*}

\subsection{Three-dimensional hopping produces a transport tensor in a real crystal geometry}
Figure~\ref{fig:bn} shows the periodic BN diffusion tensor. At $\etaperp=0.05$, the ensemble averages are $D_1=5.09\times10^{-8}$, $D_2=3.87\times10^{-8}$, and $D_3=6.45\times10^{-9}$ m$^2$ s$^{-1}$. The principal-axis anisotropy $(D_1+D_2)/(2D_3)$ is $6.95\pm0.06$ in standard error over 12 realizations, while projections on crystallographic plane normals give $7.74\pm0.07$. The absolute diffusivities inherit the chosen attempt frequency; the anisotropy ratios and principal directions do not.

\begin{figure*}[t]\centering
\includegraphics[width=0.98\textwidth]{../reproduction/generated/figure3_bn_tensor.png}
\caption{\textbf{Tensorial diffusion in an atomistic BN crystal.} Principal diffusivities and anisotropy evolve continuously with transverse hopping strength, demonstrating a robust slow out-of-plane mode at weak coupling.}\label{fig:bn}\end{figure*}

\subsection{Chemical contrast selects the dynamically active backbone}
Figure~\ref{fig:chemistry} introduces a minimal chemical landscape on the W$_2$O$_6$ graph. Across the full map, eventual passage remains unity, yet the trajectory observables change strongly. At $w_{\mathrm{homo}}=0.03$, increasing $\Delta\epsilon$ from 0 to 0.25 eV changes the reduced first-passage time from 13.99 to 167.7 and reduces the W-site visit fraction from 47.4\% to 5.2\%. The response saturates once the contrast greatly exceeds $k_BT$, after which the remaining kinetic variation is governed primarily by pair weights.

\begin{figure*}[t]\centering
\includegraphics[width=0.98\textwidth]{../reproduction/generated/figure4_chemistry.png}
\caption{\textbf{Chemical heterogeneity reshapes kinetics and path composition in W$_2$O$_6$.} Reduced first-passage times and species-resolved trajectory participation reveal strong backbone selection despite a saturated eventual passage fraction.}\label{fig:chemistry}\end{figure*}

\subsection{A reduced description of dimensionality, disorder, and chemistry}
The calculations are summarized by three independent controls. Vacancy fraction changes topology; $\etaperp$ changes the availability and kinetic cost of excursions through the third dimension; chemical contrast and pair weights change the energetic accessibility of otherwise connected channels. Their signatures differ: reachability responds to topology, first-passage time and tortuosity to kinetic bottlenecks, and species-resolved trajectory statistics to chemical selection.

The approach applies to regimes in which localized-state hopping is a meaningful coarse-grained description. It does not replace coherent quantum transport for clean low-temperature crystals, nor does it predict absolute conductivity without carrier density, device geometry, and a calibrated attempt frequency. Where first-principles information is available, site-resolved energies and edge-resolved transfer weights can be supplied without changing the trajectory algorithm.

\section{Conclusions}
We have extended a trajectory-resolved hopping model from planar atomic networks to three-dimensional and chemically heterogeneous structures while separating graph topology, kinetic accessibility, and local energetics. Weak transverse connections can restore source--drain reachability in a vacancy-fragmented network even when those links are dynamically slow; once reconnected, stronger transverse hopping reduces first-passage times. Periodic BN simulations produce a robust slow out-of-plane diffusion mode, and the chemistry-aware W$_2$O$_6$ test shows that local energetic contrast can redirect the active trajectory backbone without changing eventual passage. The framework therefore provides a practical bridge between crystal structure, disorder, and mesoscopic hopping transport while retaining a simple CIF-plus-parameter-file workflow.

\section*{Data availability}
The atomic structures, parameter files, processed data used to generate the figures, and the Python implementation required to reproduce the calculations are provided with the public code repository accompanying this work.

\section*{Acknowledgments}
The authors acknowledge computational resources from LCCMat at the University of Bras\'ilia. Funding information should be inserted here before final submission.

\bibliography{bibliography}
\end{document}
